\documentclass[aspectratio=169,11pt]{beamer}
\usetheme{default}
\setbeamertemplate{navigation symbols}{}
\definecolor{TextJEPAGreen}{HTML}{1F6B4F}
\setbeamercolor{structure}{fg=TextJEPAGreen}
\setbeamercolor{frametitle}{fg=white,bg=TextJEPAGreen}

% Use this only after REPORT.md passes validation. A slide is a teaching aid,
% not a replacement for the self-contained report. Write complete sentences,
% define every term, and use additional slides instead of shrinking content.

\title{A descriptive result title}
\subtitle{A beginner-first explanation of one TextJEPA decision}
\author{TextJEPA research loop}
\date{\today}

\begin{document}
\frame{\titlepage}

\begin{frame}{The whole result in ordinary language}
  \begin{itemize}
    \item Explain the problem with a concrete everyday example.
    \item State what was tested without unexplained abbreviations.
    \item State what happened and how uncertain that answer is.
    \item State why the result changes, or does not change, the direction.
  \end{itemize}
\end{frame}

\begin{frame}{What is a language model, and what are we asking it to learn?}
  \begin{itemize}
    \item Define a neural network as a program that adjusts many numerical knobs from examples.
    \item Define a language model before introducing JEPA or hidden representations.
    \item Explain a small token step and a larger phrase or sentence step with one example.
  \end{itemize}
\end{frame}

\begin{frame}{The idea as a picture}
  \centering
  % \includegraphics[width=.85\linewidth]{figures/main-idea.pdf}
  \vspace{2cm}
  \textbf{Place one intuitive mechanism figure here.}

  \vspace{1em}
  \small What to notice: state the visual pattern. Why it matters: connect it to the decision.
\end{frame}

\begin{frame}{What we compared, and why the comparison is fair}
  \begin{itemize}
    \item Explain what information every system receives.
    \item Explain the matched control and negative control.
    \item Explain tuning, seeds, compute, exclusions, and leakage checks.
    \item Say which apparent success would be invalid and why.
  \end{itemize}
\end{frame}

\begin{frame}{What happened}
  \centering
  \begin{tabular}{p{.30\linewidth}p{.22\linewidth}p{.38\linewidth}}
    \textbf{System} & \textbf{Result} & \textbf{Meaning in ordinary language}\\\hline
    Main idea & -- & Explain the number rather than only printing it.\\
    Fair control & -- & Explain what difference remains.\\
    Negative control & -- & Explain whether the measurement behaves sensibly.
  \end{tabular}
\end{frame}

\begin{frame}{Technical details for reproducibility}
  \begin{itemize}
    \item Define the inputs, targets, architecture, and losses.
    \item Give dataset identity, scale, seeds, optimizer, and stopping rule.
    \item Define each primary metric in words before its mathematical form.
    \item Link the exact commit, manifests, summaries, and full report.
  \end{itemize}
\end{frame}

\begin{frame}{What we can conclude}
  \begin{itemize}
    \item Separate direct observations from interpretations.
    \item State the scope: dataset, scale, model, and evaluation setting.
    \item Explain exactly which decision the evidence supports.
  \end{itemize}
\end{frame}

\begin{frame}{What we cannot conclude}
  \begin{itemize}
    \item List remaining alternative explanations and failed validity gates.
    \item Say whether the evidence supports hierarchy, generation, or only a diagnostic.
    \item State what would falsify the preferred interpretation.
  \end{itemize}
\end{frame}

\begin{frame}{What happens next, and where you should steer}
  \begin{itemize}
    \item State the smallest next question and its expected result patterns.
    \item State the compute cost and scale gate.
    \item Ask two concrete questions through which the human can change priorities.
  \end{itemize}
\end{frame}

\end{document}

